\chapter{Vectorial integration}

In this chapter, if F denotes a Hausdorff locally convex vector space (over R or C), we denote by \(F'\) its dual, by \(F''\) its bidual, and by \(F'^*\) the algebraic dual of \(F'\) (the space of all linear forms on \(F'\) ); \(F''\) is a linear subspace of \(F'^*\) , and F may be identified (as a vector space without topology) with a linear subspace of \(F''\) . We denote by \(F_\sigma\) the vector space F equipped with the weakened topology \(\sigma(F, F')\) ; the qualifiers `weak' and `weakly' refer to this topology.

In this chapter, T denotes a locally compact space, \(\mathcal{K}_{\mathbf{R}}(\mathrm{T})\) or \(\mathcal{K}(\mathrm{T})\) (resp. \(\mathcal{K}_{\mathbf{C}}(\mathrm{T})\) ) the vector space of real (resp. complex) functions on T, continuous and with compact support; for every subset A of T, \(\mathcal{K}(\mathrm{T},\mathrm{A})\) (resp. \(\mathcal{K}_{\mathbf{C}}(\mathrm{T},\mathrm{A})\) ) denotes the subspace of \(\mathcal{K}(\mathrm{T})\) (resp. \(\mathcal{K}_{\mathbf{C}}(\mathrm{T})\) ) formed by the functions whose support is contained in A. Absent express mention to the contrary, the space \(\mathcal{K}(\mathrm{T})\) (resp. \(\mathcal{K}_{\mathbf{C}}(\mathrm{T})\) ) will be equipped with the direct limit topology of the topologies of uniform convergence on each of the subspaces \(\mathcal{K}(\mathrm{T},\mathrm{K})\) (resp. \(\mathcal{K}_{\mathbf{C}}(\mathrm{T},\mathrm{K})\) ), where K runs over the set of compact subsets of T.

We recall that this topology is finer than the topology of uniform convergence, hence is Hausdorff; it induces on each \(\mathcal{K}(\mathrm{T},\mathrm{K})\) (resp. \(\mathcal{K}_{\mathbf{C}}(\mathrm{T},\mathrm{K})\) ) the topology of uniform convergence (TVS, II, §4, No.~4, Remark). The space \(\mathcal{K}_{\mathbf{C}}(\mathrm{T})\) may be identified with the space obtained from \(\mathcal{K}(\mathrm{T})\) by extension of the scalars from R to C. To say that a linear form on \(\mathcal{K}(\mathrm{T})\) is a measure means, by definition, that it is continuous (Ch. III, §1, No.~3, Def. 2).

